\subsection{REGULATED FUNCTIONS}

DEFINITION 2. One says that a map f of an interval \(I \subset R\) into a set E is a step function if there is a partition of I into a finite number of intervals \(J_{k}\) such that f is constant on each of the \(J_{k}\) .

Let \((a_{i})_{0\leqslant i\leqslant n}\) be the strictly increasing sequence formed by the distinct endpoints of the \(\mathrm{J}_k\) ; since the \(\mathrm{J}_k\) are pairwise disjoint each of them either reduces to a singleton \(a_{i}\) or is an interval with two consecutive points \(a_{i}\) , \(a_{i + 1}\) as endpoints; moreover, since I is the union of the \(\mathrm{J}_k\) , the point \(a_0\) is the left-hand endpoint, and \(a_{n}\) is the right-hand endpoint of I. Every step function on I can thus be characterised as a function which is constant on each of the open intervals \(]a_{i}\) , \(a_{i+1}[ (0 \leqslant i \leqslant n-1)\) , where \((a_{i})_{0 \leqslant i \leqslant n}\) is a strictly increasing sequence of points of I with \(a_{0}\) being the left-hand endpoint and \(a_{n}\) the right-hand endpoint of I.

PROPOSITION 2. The set of step functions defined on I, with values in a vector space E over R, is a vector subspace E of the vector space \(\mathcal{F}(\mathrm{I};\mathrm{E})\) of all maps of I into E.

Indeed, let \(f\) and \(g\) be two step functions, and \((\mathrm{A}_i)\) and \((\mathrm{B}_j)\) two partitions of I into a finite number of intervals such that \(f\) (resp. \(g\) ) is constant on each of the \(\mathrm{A}_i\) (resp. \(\mathrm{B}_j\) ); whatever the real numbers \(\lambda, \mu\) , it is clear that \(\lambda f + \mu g\) is constant on each of the nonempty intervals \(\mathrm{A}_i \cap \mathrm{B}_j\) , and that these intervals form a partition of I.

COROLLARY. The vector subspace E is generated by the characteristic functions of intervals.

Now let us consider the case where E is a normed space over R; it is then immediate that the characteristic function of an interval J with endpoints a, b (a \textless{} b) admits a primitive, namely the function equal to a for \(x \leqslant a\) , to x for \(a \leqslant x \leqslant b\) , and to b for \(x \geqslant b\) . The cor. to prop. 2 thus shows that every step function with values in E admits a primitive.

We can now apply the method set out in \(n^{\circ}\) 2.

DEFINITION 3. One says that a vector function, defined on an interval I, with values in a complete normed space E over R, is a regulated function, if it is the uniform limit of step functions on every compact subset of I.

In other words, the regulated functions are the elements of the closure in \(\mathcal{F}_{c}(\mathrm{I};\mathrm{E})\) of the subspace E of step functions; \(\overline{E}\) is a vector subspace of \(\mathcal{F}_{c}(\mathrm{I};\mathrm{E})\) and since \(\mathcal{F}_{c}(\mathrm{I};\mathrm{E})\) is complete, so is \(\overline{E}\) ; in other words, if a function is the uniform limit of regulated functions on every compact subset of I, then it is regulated on I. For f to be regulated on an interval I it is necessary and sufficient that its restriction to every compact interval contained in I be regulated.

Cor. 1 to II, p.~53 shows:

THEOREM 2. Every regulated function on an interval I admits a primitive on I.

We shall transform def. 3 of II, p.~4 to another equivalent one:

THEOREM 3. For a vector function f defined on an interval I, with values in a complete normed space E over R, to be regulated, it is necessary and sufficient that it have a right limit and a left limit at every interior point of I, and a right limit at the left-hand endpoint of I and a left limit at the right-hand endpoint of I, when these points belong to I. The set of points of discontinuity of f in I is thus countable.

Since every interval is a countable union of compact intervals, one can restrict oneself to proving th. 3 when I is compact, say I = {[}a, b{]}.

\(1^{\circ}\) The condition is necessary. Suppose that \(\mathbf{f}\) is regulated and let \(x\) be a point of I different from \(b\) . By hypothesis, for every \(\varepsilon > 0\) there is a step function \(\mathbf{g}\) such that \(\| \mathbf{f}(z) - \mathbf{g}(z) \| \leqslant \varepsilon\) for every \(z \in \mathrm{I}\) ; since \(\mathbf{g}\) has a right limit at the point \(x\) there exists a \(y\) such that \(x < y \leqslant b\) and such that, for every pair of points \(z\) , \(z'\) in the interval \([x, y]\) one has \(\| \mathbf{g}(z) - \mathbf{g}(z') \| \leqslant \varepsilon\) , and in consequence \(\| \mathbf{f}(z) - \mathbf{f}(z') \| \leqslant 3\varepsilon\) ; this proves (by Cauchy's criterion) that \(\mathbf{f}\) has a right limit at the point \(x\) . In the same way one proves that \(\mathbf{f}\) has a left limit at every point of I different from \(a\) .

\(2^{\circ}\) The condition is sufficient. Suppose it is satisfied; for each \(x \in I\) there is an open interval \(\mathrm{V}_x = ]c_x\) , \(d_x[\) containing \(x\) and such that on the intersection of \(I\) with each of the open intervals \(]c_x\) , \(x[i, ]x\) , \(d_x[\) (when the intersection is not empty) the oscillation of \(\mathbf{f}\) is \(\leqslant \varepsilon\) . Since \(I\) is compact there is a finite number of points \(x_i\) in \(I\) such that the \(\mathrm{V}_{x_i}\) form a cover of \(I\) ; let \((a_k)_{0 \leqslant k \leqslant n}\) be the sequence obtained by arranging in increasing order the points of the finite set formed by the points \(a\) , \(b\) and those points \(x_i\) , \(c_{x_i}\) and \(d_{x_i}\) which belong to \(I\) ; each of the intervals \(]a_k\) , \(a_{k+1}[\) ( \(0 \leqslant k \leqslant n-1\) ) being contained in an interval \(]c_{x_i}\) , \(x_i[\) or \(]x_i\) , \(d_{x_i}[\) , the oscillation of \(\mathbf{f}\) there is \(\leqslant \varepsilon\) ; let \(\mathbf{c}_k\) be one of the values of \(\mathbf{f}\) on \(]a_k\) , \(a_{k+1}[\) ; on putting \(\mathbf{g}(a_k) = \mathbf{f}(a_k)\) for \(0 \leqslant k \leqslant n\) , and \(\mathbf{g}(x) = \mathbf{c}_k\) for all \(x \in ]a_k\) , \(a_{k+1}[\) ( \(0 \leqslant k \leqslant n-1\) ), one defines a step function \(\mathbf{g}\) such that \(\| \mathbf{f}(z) - \mathbf{g}(z) \| \leqslant \varepsilon\) on \(I\) ; so \(\mathbf{f}\) is regulated on \(I\) .

Finally, let us show that if \(\mathbf{f}\) is regulated on I then its set of points of discontinuity is countable. For every \(n > 0\) there exists a step function \(\mathbf{g}_n\) such that \(\| \mathbf{f}(x) - \mathbf{g}_n(x)\| \leqslant 1 / n\) on I; since the sequence \((\mathbf{g}_n)\) converges uniformly to \(\mathbf{f}\) on I, we see that \(\mathbf{f}\) is continuous at every point where the \(\mathbf{g}_n\) are all continuous (Gen.~Top., X, p.~281, cor. 1); but since \(\mathbf{g}_n\) is continuous except at the points of a finite set \(\mathrm{H}_n\) it follows that \(\mathbf{f}\) is continuous at the points of the complement of the set \(\mathrm{H} = \bigcup_{n}\mathrm{H}_{n}\) , which is countable.

COROLLARY 1. Let \(\mathbf{f}\) be a regulated function on I; at every point of I, apart from the right-hand endpoint (resp. the left-hand endpoint) of I, every primitive of \(\mathbf{f}\) has a right derivative equal to \(\mathbf{f}(x+)\) (resp. a left derivative equal to \(\mathbf{f}(x-)\) ); in particular, at every point \(x\) where \(\mathbf{f}\) is continuous, \(\mathbf{f}(x)\) is the derivative of one of its primitives.

This is an immediate consequence of th. 3 of II, p.~54 and of prop. 6 of I, p.~18.

COROLLARY 2. Let \(\mathbf{f}_i\) ( \(1 \leqslant i \leqslant n\) ) be \(n\) regulated functions on an interval I, each \(\mathbf{f}_i\) having its values in a complete normed space \(\mathrm{E}_i\) over \(\mathbf{R}\) ( \(1 \leqslant i \leqslant n\) ). If \(\mathbf{g}\) is a continuous map of the subspace \(\prod_{i=1}^{n} \overline{\mathbf{f}_i(\mathrm{I})}\) of \(\prod_{i=1}^{n} \mathrm{E}_i\) into a complete normed space F over \(\mathbf{R}\) , then the composite function \(x \mapsto \mathbf{g}(\mathbf{f}_1(x), \mathbf{f}_2(x), \ldots, \mathbf{f}_n(x))\) is regulated on I.

Indeed, it clearly satisfies the conditions of th. 3 of II, p.~54.

Thus one sees that if \(\mathbf{f}\) is a regulated vector function on I, then the real function \(x\mapsto \| \mathbf{f}(x)\|\) is also regulated. Further, the real regulated functions on I form a ring;

moreover, if \(f\) and \(g\) are two real regulated functions, then \(\sup(f, g)\) and \(\inf(f, g)\) are regulated.

Remark. 1) If f is a real regulated function on I, and g is a regulated vector function on an interval containing \(f(\mathrm{I})\) , then the composite function \(g \circ f\) is not necessarily regulated (cf.~II, p.~79, exerc. 4).

Two particular cases of th. 3 of II, p.~54 are especially important:

PROPOSITION 3. Every continuous vector function on an interval \(I \subset R\) taking its values in a complete normed space E over R is regulated, and admits a primitive on I, of which it is the derivative at every point.

Remarks. 2) To show that a continuous function admits a primitive, one can use the fact that every polynomial function of a real variable (with coefficients in E) admits a primitive; since from the theorem of Weierstrass (Gen.~Top., X, p.~313, prop. 3) every continuous function is the uniform limit of polynomials on every compact interval, th. 1 of II, p.~52 shows that every continuous function admits a primitive.

\begin{enumerate}
\def\labelenumi{\arabic{enumi})}
\setcounter{enumi}{2}
\tightlist
\item
  The principle of the preceding remark extends without significant modification to vector functions of a complex variable taking values in a complete normed space over \(\mathbf{C}\) . If \(\mathrm{U}\) is an open set in \(\mathbf{C}\) , homeomorphic to \(\mathbf{C}\) , a primitive of such a vector function \(\mathbf{f}\) defined on \(\mathrm{U}\) is by definition a continuous function on \(\mathrm{U}\) , having derivative equal to \(\mathbf{f}\) at every point of \(\mathrm{U}\) . With this definition, th. 1 of II, p.~52 extends without modification (one proves, using the connectedness of \(\mathrm{U}\) , that \((\mathbf{g}_{\alpha})\) is uniformly convergent with respect to \(\mathcal{F}\) on a neighbourhood of each point of \(\mathrm{U}\) , from which it follows that \((\mathbf{g}_{\alpha})\) is uniformly convergent with respect to \(\mathcal{F}\) on every compact subset of \(\mathrm{U}\) ; the proof is completed using prop. 4 of I, p.~18). Consequently, every function which is a uniform limit of polynomials on every compact subset of \(\mathrm{U}\) , admits a primitive on \(\mathrm{U}\) ; these functions are no other than the functions called holomorphic on \(\mathrm{U}\) , which we shall study further in detail in a later Book.
\end{enumerate}

PROPOSITION 4. Every monotone real function f on an interval \(I \subset R\) is regulated, and every primitive of f is convex on I.

Indeed, f satisfies the criterion of th. 3 of Gen.~Top., IV, p.~350, prop. 4; the second part of the proposition follows from cor. 1, from II, p.~55, and from prop. 5 of I, p.~27.

Remark. 4) One must not think that the regulated functions on an interval I are the only functions having a primitive on I (cf.~II, p.~80, exercises 7 and 8).

